% META: {"id": "TNSI-ARCHITECTURES-MATERIELLES-SY-RE-C02B", "chapitre": "TNSI-ARCHITECTURES-MATERIELLES-SY", "type_objet": "remediation", "capacites_codes": ["C02B"], "status": "generated"}

%---------------------------------------------------------------
% FICHE DE REMEDIATION — C02B : concurrence vs parallelisme
%---------------------------------------------------------------

\begin{exercice}{TNSI-ARCH-RE-C02B-EX1}{1}{10}
Un ordinateur possède un processeur à un seul cœur, et pourtant l'utilisateur peut écouter de la musique tout en tapant du texte dans un traitement de texte, sans interruption perceptible. Ces deux tâches s'exécutent-elles en parallélisme réel ou en concurrence ? Expliquer.
\end{exercice}

% BEGIN-VERIFY
% def tourniquet(processus, tranche, duree_totale):
%     ordre_execution = []
%     file = list(processus)
%     temps_restant = dict(duree_totale)
%     while file:
%         p = file.pop(0)
%         execute = min(tranche, temps_restant[p])
%         ordre_execution.append((p, execute))
%         temps_restant[p] -= execute
%         if temps_restant[p] > 0:
%             file.append(p)
%     return ordre_execution
% resultat = tourniquet(["musique", "texte"], 1, {"musique": 2, "texte": 2})
% assert resultat == [("musique", 1), ("texte", 1), ("musique", 1), ("texte", 1)]
% END-VERIFY

\begin{corrige}{TNSI-ARCH-RE-C02B-EX1}
Ces deux tâches s'exécutent en \textbf{concurrence}, pas en parallélisme réel : un seul cœur ne peut exécuter qu'une instruction à la fois. L'ordonnanceur alterne très rapidement entre les deux processus (tranches de temps de quelques millisecondes), ce qui donne à l'utilisateur l'\textbf{illusion} d'une exécution simultanée, alors qu'en réalité chaque processus n'avance que pendant sa tranche de temps.
\end{corrige}
